% Part: first-order-logic
% Chapter: syntax-and-semantics
% Section: main-operator

\documentclass[../../../include/open-logic-section]{subfiles}

\begin{document}

\olfileid[hi]{fol}{syn}{mai}

\olsection{किसी सूत्र का \printtoken{S}{main operator}}

\begin{explain}
किसी !!a{formula}~$!A$ के निर्माण में सबसे अंत में प्रयुक्त संक्रियक की
चर्चा करना प्रायः उपयोगी होता है। इस संक्रियक को~$!A$ का \emph{मुख्य
संक्रियक} कहते हैं। सहज रूप से यह $!A$ का ``सबसे बाहरी'' संक्रियक है।
उदाहरणतः $\lnot !A$ का मुख्य संक्रियक $\lnot$ है, $(!A \lor !B)$ का मुख्य
संक्रियक $\lor$ है, इत्यादि।
\end{explain}


\begin{defn}[!!^{main operator}]
\ollabel{def:main-op}
किसी !!a{formula}~$!A$ का \emph{!!{main operator}} निम्न प्रकार परिभाषित है:
\begin{enumerate}
\item \indcase*{!A}{!A}{$\indfrm$ का कोई !!{main operator} नहीं है।}

\tagitem{prvNot}{\indcase{!A}{\lnot !B}{$\indfrm$ का !!{main operator}
  ~ $\lnot$ है।}}{}

\tagitem{prvAnd}{\indcase{!A}{(!B \land !C)}{$\indfrm$ का
  !!{main operator}~$\land$ है।}}{}

\tagitem{prvOr}{\indcase{!A}{(!B \lor !C)}{$\indfrm$ का
  !!{main operator}~$\lor$ है।}}{}

\tagitem{prvIf}{\indcase{!A}{(!B \lif !C)}{$\indfrm$ का
  !!{main operator}~$\lif$ है।}}{}

\tagitem{prvIff}{\indcase{!A}{(!B \liff !C)}{$\indfrm$ का
  !!{main operator}~$\liff$ है।}}{}

\tagitem{prvAll}{\indcase{!A}{\lforall[x][!B]}{$\indfrm$ का
  !!{main operator}~$\lforall$ है।}}{}

\tagitem{prvEx}{\indcase{!A}{\lexists[x][!B]}{$\indfrm$ का
  !!{main operator}~$\lexists$ है।}}{}
\end{enumerate}
\end{defn}

हर स्थिति में हमारा आशय सूत्र में !!{main operator} की दर्शाई गई विशिष्ट
\emph{उपस्थिति} से है। उदाहरणतः सूत्र $((!D \lif !E) \lif (!E \lif !D))$,
$(!B \lif !C)$ के रूप में है, जहाँ $!B$, $(!D \lif !E)$ है और $!C$,
$(!E \lif !D)$ है; इसलिए $\lif$ की दूसरी उपस्थिति !!{main operator} है।

\begin{explain}
यह उस फलन की \emph{पुनरावर्ती} परिभाषा है जो सभी अपरमाण्विक !!{formula}s
को उनके !!{main operator} की उपस्थिति पर प्रतिचित्रित करता है। !!{formula}s
जिस ढंग से आगमनात्मक रूप से परिभाषित हैं, उसके कारण प्रत्येक
!!{formula}~$!A$, \olref{def:main-op} के मामलों में से किसी एक को संतुष्ट
करता है। इससे प्रत्येक अपरमाण्विक !!{formula}~$!A$ के लिए !!a{main
  operator} का अस्तित्व सुनिश्चित होता है। चूँकि प्रत्येक !!{formula} इन
शर्तों में से केवल एक को संतुष्ट करता है, और हर स्थिति में जिन छोटे
!!{formula}s से $!A$ बना है वे अद्वितीय रूप से निर्धारित हैं, इसलिए
!!{main
  operator} की~$!A$ में उपस्थिति अद्वितीय है और हमने सचमुच एक फलन परिभाषित
किया है।
\end{explain}

किसी !!{formula} के !!{main operator} का प्रतीक कौन-सा है, इसके अनुसार हम
उसे \olref{tab:main-op} में दिए नाम से बुलाते हैं।
\iftag{defNot,defOr,defAnd,defIf,defIff,defTrue,defFalse,defEx,defAll}
{ किंतु याद रहे कि परिभाषित संक्रियक औपचारिक रूप से !!{formula}s में नहीं
आते। वे केवल संक्षेप हैं, इसलिए औपचारिक रूप से किसी सूत्र के मुख्य संक्रियक
नहीं हो सकते। अतः सभी !!{formula}s के बारे में प्रमाणों में उन्हें अलग से
संभालने की आवश्यकता नहीं है।}

\begin{table}[!h]
\centering
\begin{tabular}{c | c | c}
!!^{main operator} & !!{formula} का प्रकार & उदाहरण\\
\hline
कोई नहीं & परमाण्विक (!!{formula}) &
\iftag{prvFalse}{$\lfalse$,}{}
\iftag{prvTrue}{$\ltrue$,}{}
$\Atom{R}{t_1, \dots, t_n}$,
$\eq[t_1][t_2]$\\
$\lnot$ & निषेध & $\lnot !A$ \\
$\land$ & संयोजन & $(!A \land !B$) \\
$\lor$ & वियोजन & $(!A \lor !B$) \\
$\lif$ & !!{conditional} & $(!A \lif !B$) \\
$\liff$ & !!{biconditional} & $(!A \liff !B)$ \\
$\lforall[][]$ & सार्वत्रिक (!!{formula})& $\lforall[x][!A]$ \\
$\lexists[][]$ & अस्तित्वपरक (!!{formula})& $\lexists[x][!A]$
\end{tabular}
\caption{!!{formula}s के मुख्य संक्रियक और नाम}
\ollabel{tab:main-op}
\end{table}

\end{document}
